problem:  
Let \((X,d,\mathfrak{m})\) be a compact, essentially non‑branching metric measure space satisfying the curvature–dimension condition \(\mathrm{CD}(K,N)\) for some \(K\in\mathbb{R}\), \(N>1\).  
For each pair of absolutely continuous measures \(\mu_0,\mu_1\in\mathcal P_2(X)\) with bounded densities, denote by \((\mu_t^\varepsilon)_{t\in[0,1]}\) the entropic interpolation between them with regularization parameter \(\varepsilon>0\), and by \((\mu_t)\) the limiting Wasserstein geodesic as \(\varepsilon\to 0\).  

Define the **metric Fisher–information deviation functional**  
\[
\mathcal{I}_\varepsilon(\mu_0,\mu_1)
:= \int_0^1\!\!\Big(
I(\mu_t^\varepsilon) - I(\mu_t)
\Big)\,dt,
\quad 
I(\nu):=\int_X |\nabla^- \log\rho_\nu|^2\,d\nu,
\]
where \(I(\nu)\) denotes the metric Fisher information defined in terms of the local slope \(|\nabla^- \log\rho_\nu|\) of the entropy, which is always available in CD spaces.  

Assume that for every \(\mu_0,\mu_1\) as above, the limit  
\[
Q(\mu_0,\mu_1):=\lim_{\varepsilon\to0}\varepsilon^{-2}\mathcal{I}_\varepsilon(\mu_0,\mu_1)
\]
exists and defines a finite, symmetric, quadratic functional on absolutely continuous measures.  

**Question:**  
Does the existence of such a quadratic limit \(Q\) for all pairs \((\mu_0,\mu_1)\) imply that \((X,d,\mathfrak{m})\) has a linear heat flow, i.e. that its Cheeger energy is quadratic and it is in fact an \(\mathrm{RCD}(K,N)\) space?  
Equivalently, can the **second‑order Fisher‑information asymptotics** of the entropic interpolation serve as a quantitative analytic test distinguishing \(\mathrm{CD}(K,N)\) from \(\mathrm{RCD}(K,N)\) spaces?  

Why is it a "good" problem:  
This formulation avoids the circularity of requiring Hessians, relying only on first‑order quantities (entropy gradients and Fisher information) that are well-defined for all CD spaces. Yet the question still fundamentally probes the emergence of *second‑order geometric structure* (the quadraticity of Cheeger energy and the tensorial Ricci notion) from entropic transport asymptotics.  

1. **Conceptual bridge:**  
   The problem connects *stochastic smoothing via the Schrödinger problem* with the *first‑ and second‑order calculus of metric measure spaces*. It asks whether higher‑order geometric structure can be inferred purely from the decay rate of the Fisher‑information gap, forging a deep link between probabilistic regularization and differential geometry.  

2. **Mathematical soundness:**  
   Every term in the statement—local slope, Fisher information, entropic interpolation—exists in any CD space. Thus, the problem is logically well‑posed even without second‑order calculus.  

3. **Depth and simplicity:**  
   Expressed with familiar analytic objects, the problem carries profound meaning: it seeks whether a curvature–dimension space whose entropic interpolations exhibit *quadratic Fisher‑information stability* must necessarily possess a linear (Riemannian) heat flow. This blends probabilistic asymptotics, information geometry, and the structure theory of metric measure spaces.  

4. **Potential theoretical impact:**  
   A positive answer would create a new analytic characterization of RCD spaces through Schrödinger‑bridge asymptotics, paralleling how linearity of the heat flow characterizes RCD structure, while grounding it in measurable transport quantities accessible to stochastic analysis.  

Hence this version constitutes a mathematically sound and conceptually deep “good” problem uniting synthetic geometry, analysis, and probability.